% =====================================================================
%  3 Funktionale Anforderungen  (Aufgabenstellung: „Detaillierte
%  Auflistung der gewünschten Funktionen und Features. Die Anforderungen
%  sollten textuell durch User Stories beschrieben und durch ein
%  UML-Use-Case-Diagramm visualisiert werden.“)
%  Format: „Als <Rolle> möchte ich <Ziel>, damit <Nutzen>.“ (Cohn, 2004)
%  IDs (US01 …) werden in Spezifikation, Architektur und Testdokument
%  wiederverwendet (Nachverfolgbarkeit!).
% =====================================================================
\section{Funktionale Anforderungen}
\TODO{Kurze Hinführung}

\subsection{User Stories}
\TODO{User Stories nach Bereichen der App gruppieren; Priorität nach MoSCoW; je Story Akzeptanzkriterien}

\begin{longtable}{@{}L{1.1cm}L{8.2cm}L{1.3cm}L{4.6cm}@{}}
  \caption{User Stories}\label{tab:userstories}\\
  \toprule
  \textbf{ID} & \textbf{User Story} & \textbf{Prio.} & \textbf{Akzeptanzkriterien} \\
  \midrule
  \endfirsthead
  \toprule
  \textbf{ID} & \textbf{User Story} & \textbf{Prio.} & \textbf{Akzeptanzkriterien} \\
  \midrule
  \endhead
  \bottomrule
  \endfoot
  US01 & Als \TODO{Rolle} möchte ich \TODO{Ziel}, damit \TODO{Nutzen}. & Must & \TODO{} \\
  US02 & Als \TODO{Rolle} möchte ich \TODO{Ziel}, damit \TODO{Nutzen}. & Should & \TODO{} \\
\end{longtable}
\quelle{Eigene Darstellung.}

\subsection{Use-Case-Diagramm}
\TODO{UML-Use-Case-Diagramm (Akteure, Systemgrenze, Use Cases, include/extend) in draw.io erstellen: docs/diagramme/01\_UseCase/ – Export als docs/diagramme/export/usecase.pdf; kurz erläutern}
\begin{figure}[H]
  \caption{Use-Case-Diagramm}\label{abb:usecase}
  \centering
  \grafik{usecase}
  \quelle{Eigene Darstellung.}
\end{figure}
